\scanpage{159}
\mspage{75}
% archive sheet number: 158

a alors

\[
x\overset{\mathrm{def}}{=}\underset{I}{\underrightarrow{\lim}}\ \underset{\rule{.3pt}{2pt}\rule{1.2em}{.3pt}\rule{.3pt}{2pt}}{x_i}
=\underset{J}{\underrightarrow{\lim}}\ \underset{\substack{\rule{.3pt}{2pt}\rule{1.2em}{.3pt}\rule{.3pt}{2pt}\\\shortparallel\\\psi(j)}}{\unsure{\zeta|\xi}_{\unsure{j|s}}},
\]

\margin{\noindent\begin{tabular}{@{}|p{.92\linewidth}@{}}
On utilise la propriété de distributivité des $\underrightarrow{\lim}$, rel. à un foncteur Cat-cofibrant \diagram{\begin{tikzcd}[ampersand replacement=\&,row sep=small,column sep=small]
J\arrow[r]\arrow[dr,"\psi"'] \& I\\
\& M
\end{tikzcd}} explicitée dans
\end{tabular}}

où cette fois la limite est une limite d'objets de $S$, \inserted{prise dans $M$.} Il s'ensuit que $x\in\mathrm{Filt}^{\prime\,\pi}$, ce n'est pas tout à fait $x\in\mathrm{Filt}^{\pi}$, \deleted{Mais sauf \unsure{er}} [\unsure{ça|ce} implique $x\in\mathrm{Filt}^{\pi^{\pi_0}}$, \deleted{\unsure{et}} donc $x\in\mathrm{Filt}^{\pi}$ si $\pi=\pi^{\pi_0}$, p. ex. $\pi=2^c$ avec $c\geq\pi_0$]

Mais c'est le moment d'utiliser le fait que les objets de $S$ sont accessibles, et on prendra $\pi_0$ tel que les objets de $S$ soient $\pi_0$-accessibles, et cette fois on prendra quand-même $I$ \deleted{filtrante} et grande devant $\pi_0$ (pour nous borner à ce qui est exigé dans la condition b) de loc. cit.) On doit prouver que $\mathrm{Card}(\mathrm{Fl}\ S_{/x})\leq\pi$ (ce qui signifie $x\in\mathrm{Filt}^{\pi}$). Mais comme les objets de $S$ sont $\pi_0$-accessibles et $I$ grand devant $\unsure{k_0|\pi_0}$, on voit aussitôt que

\[
S_{/x}\simeq\underset{i}{\underrightarrow{\lim}}\ S_{/x_i}
\]

d'où

\[
\mathrm{Fl}\,S_{/x}\simeq\underset{i}{\underrightarrow{\lim}}\ \mathrm{Fl}\,S_{/x_i}
\]

donc

\[
\mathrm{Card}\,\mathrm{Fl}\,S_{/x}\leq
\underbrace{(\mathrm{Card}\ I)}_{\leq\pi}
\underbrace{\bigl(\underset{i}{\mathrm{Sup}}\ \mathrm{Card}\,\mathrm{Fl}\,S_{/x_i}\bigr)}_{\leq\pi}
\leq\pi^2=\pi
\]

\unsure{qed|cqfd}.

